\section{Población, muestra y variables}
\subsection{Población objetivo}
Estudiantes de 18 años o más matriculados durante el periodo de estudio en programas de educación superior de instituciones ubicadas en Bogotá, D. C. La definición se plantea para delimitar la propuesta; antes de aplicar el instrumento se debe precisar el periodo académico, los niveles de formación incluidos y las instituciones participantes. Cada estudiante participante que cumpla los criterios de inclusión y responda voluntariamente el cuestionario.
\subsection{Criterios de inclusión y exclusión}
Inclusión: tener 18 años o más, estar matriculado en un programa de educación superior de una institución ubicada en Bogotá y aceptar el consentimiento informado. Exclusión: no cumplir alguno de esos criterios o entregar un cuestionario tan incompleto que impida analizar las variables principales. No se recogerán nombres, documentos de identidad ni contraseñas.
\subsection{Diseño de muestra propuesto}
La opción preferible, si se obtiene un listado de estudiantes y autorización institucional, es un muestreo probabilístico estratificado por institución o nivel de formación, con selección aleatoria dentro de cada estrato. Este procedimiento requiere un marco muestral, información sobre el tamaño de la población y una definición previa de precisión y nivel de confianza. En ausencia de esos elementos no es válido afirmar un tamaño muestral representativo.
Si solo se dispone de participantes voluntarios o de contactos accesibles, se utilizaría una muestra no probabilística por conveniencia, documentando claramente esta limitación. \textcite{argibay2009muestra} discute la diferencia entre muestreo probabilístico y no probabilístico y advierte que la representatividad y la validez externa dependen del procedimiento de selección, no únicamente del número de casos.
\subsection{Tamaño de muestra}
El tamaño definitivo debe calcularse después de conocer el tamaño de la población accesible y decidir el nivel de confianza, el margen de error y la proporción esperada para el indicador principal. Si se emplea una muestra por conveniencia, el número de respuestas se definirá con base en la viabilidad y se reportará como muestra de participantes, sin atribuirle representatividad estadística. En esta fase no se asigna una cifra ficticia.
\newpage
\subsection{Variabilidad e indicadores}
\begin{table}[H]
    \caption{Variables e indicadores propuestos}
    \label{tab:variables-indicadores}
    \begin{tabularx}{\textwidth}{
        >{\raggedright\arraybackslash}p{0.23\textwidth}
        >{\raggedright\arraybackslash}X
        >{\raggedright\arraybackslash}p{0.19\textwidth}
    }
        \toprule
        \textbf{Variable} &
        \textbf{Indicador / operacionalización} &
        \textbf{Tipo de dato} \\
        \midrule

        Herramientas utilizadas &
        Categorías seleccionadas: LMS, buscadores, ofimática, videoconferencia, colaboración, IA y otras. &
        Categórica múltiple \\

        Frecuencia de uso &
        Nunca, rara vez, semanal, varias veces por semana, diariamente. &
        Ordinal \\

        Finalidad académica &
        Consultar fuentes, redactar, practicar, comunicarse, organizar tareas y crear contenidos. &
        Categórica múltiple \\

        Utilidad percibida &
        Escala de respuesta de 1 a 5, de nada útil a muy útil. &
        Ordinal \\

        Dificultades reportadas &
        Conectividad, dispositivo, costo, habilidades, accesibilidad, privacidad y otras. &
        Categórica múltiple \\

        Características académicas &
        Nivel de formación, modalidad y semestre o periodo (opcional). &
        Categórica / ordinal \\
        \bottomrule
    \end{tabularx}

    \begin{minipage}{\textwidth}
        \footnotesize
        \textit{Nota.} Las categorías son una propuesta inicial; deben revisarse mediante juicio experto y prueba piloto antes de la aplicación definitiva.
    \end{minipage}
\end{table}
